\chapterhead{53}{ORR-14}{Case Study: Investor Protection and Compliance Risks}{2}{2}

\lohead{53}{a}{Summarize important regulations designed to protect investors in
financial instruments, including MiFID, MiFID II, and Dodd-Frank.}

Investor protection regulations aim to ensure investors are informed about
financial transactions and protected from fraud.

\orsub{Key areas of focus}

\begin{lettered}
\item \textbf{Client suitability, disclosure, fiduciary duties.} Financial
products should be suitable for the investor, and all risks should be disclosed.
Advisors must act in the client's best interest.
\item \textbf{Preventing improper business practices.} Regulations aim to prevent
misrepresentation, market abuse, and conflicts of interest.
\end{lettered}

\orsub{Major regulations}

\noindent\begin{minipage}{\textwidth}
\renewcommand{\arraystretch}{1.32}
\begin{tabularx}{\textwidth}{@{}L{0.26\textwidth}Y@{}}
\rowcolor{navy} \thd{Regulation} & \thd{What it does}\\
{\tabfont \textbf{a) MiFID I} (2007)} &
{\tabfont Established in the EU to describe proper business conduct, regulatory
disclosure, and reporting to prevent market abuse.}\\
\rowcolor{paleblue}
{\tabfont \textbf{b) MiFID II} (2014) \textbf{and MiFIR}} &
{\tabfont Enhanced regulations post 2007--2009 financial crisis, focusing on
trader and advisor pay, conflicts of interest, fair communication, independent
advice, sales and product governance, and best execution with counterparties.}\\
{\tabfont \textbf{c) US Dodd-Frank Act}} &
{\tabfont Protects whistleblowers, regulates OTC derivatives, and limits risky
bank activities. Created the Consumer Financial Protection Bureau (CFPB). The
\textbf{Volcker Rule} limits commercial banks from engaging in speculative and
proprietary trading.}\\
\end{tabularx}
\end{minipage}
\orfigcap{The three regulatory regimes named in this reading.}

\orsub{Compliance risk management}

\begin{itemize}
\item Focuses on preventing unintentional (human error, ineffective policies) and
intentional (fraud) compliance breaches.
\item Breaches can result from poor employee engagement, weak ethics culture, or
insufficient resources.
\item Efficient regulation includes supervising employees and trades, enhancing
operations, providing adequate training, and fostering a strong ethics culture.
\end{itemize}

% ---------------------------------------------------------------------
\lohead{53}{b}{Describe and provide lessons learned from the case studies
involving violations of investor protection or compliance regulations.}

\orsub{Investor protection case studies}

Regulatory fines punish investment banks for violating investor protection rules.
These fines aim to:

\begin{checks}
\item Punish past misconduct.
\item Deter future violations by the same firm.
\item Encourage better compliance practices across the industry.
\end{checks}

\noindent\begin{minipage}{\textwidth}
\renewcommand{\arraystretch}{1.35}
\begin{tabularx}{\textwidth}{@{}L{0.24\textwidth}L{0.17\textwidth}Y@{}}
\rowcolor{navy} \thd{Firm and year} & \thd{Fine} & \thd{Violation}\\
{\tabfont \textbf{UBS} (2008)} & {\tabfont \$11.15 billion} &
{\tabfont Selling risky securities as safe.}\\
\rowcolor{paleblue}
{\tabfont \textbf{JPMorgan} (2020)} & {\tabfont \$920 million} &
{\tabfont Manipulating markets using ``spoofing'' tactics.}\\
{\tabfont \textbf{Deutsche Bank} (2022)} & {\tabfont \$2 million} &
{\tabfont Failing to get best execution for client trades; did not prioritise
best price or speed.}\\
\end{tabularx}
\end{minipage}
\orfigcap{Three enforcement actions, and the scale of the penalties.}

\begin{orexambox}{The takeaway}
Fines can be massive for serious violations like misrepresentation. Even smaller
fines can be significant for non-manipulative issues like order execution.
\end{orexambox}
